\section{Effective nerve comparisons}\label{sec:nerve}

A cover may make a complicated space easier to describe. Its nerve remembers
which subfamilies meet and forgets the shapes of their intersections. When does
it preserve topology, and who supplies the local filling data? The repository's
topology report explicitly leaves an effective nerve theorem with produced
contractions as a research problem.\footnote{Under
\texttt{docs/proofs/topology}, see
\texttt{ipra\_topology\_framework\_2026-11-15}, section 12,
and the homology entry of its difficulty ledger.}
For a cover by simplices we construct a rational
piecewise-linear homotopy equivalence, including its homotopies. A cover by
arbitrary subcomplexes admits a different test based on linear algebra. This
one produces local contractions or an obstruction, and in the successful
case gives a natural homology isomorphism. For covers with collapsible intersections we then produce geometric contractions and a homotopy inverse, including on the three-arc cover of a hexagon.

\subsection{A homotopy equivalence one can evaluate}

A simplicial complex $K$ on $v$ vertices is a downward-closed list of
nonempty subsets, with the empty set treated separately. Code a subset by
its $v$-bit mask. Assume first that the listed cover consists of simplices
$S_1,\ldots,S_r$ and all their nonempty faces. Every simplex of $K$ must
lie in one of the $S_i$. The nerve $N$ contains a nonempty index set $I$
exactly when $\bigcap_{i\in I}S_i$ is nonempty. For example, $S_i$ may
be the maximal faces of any finite simplicial complex, found by a scan.
For a finite poset $P$, its order complex has the nonempty strict chains
of $P$ as simplices. The order complex of the nonempty face poset $P(K)$
is the barycentric subdivision $\operatorname{sd}K$. A rational point of a
realization is a nonnegative rational weight vector of sum one whose
positive support is a simplex. A named point is a regular Cauchy name of
such vectors in the sup norm. This lets consecutive samples lie on
different simplices meeting along a boundary.

\begin{theorem}\label{thm:nerve-pl}
There are rational piecewise-linear maps between the realizations of $K$
and $N$ whose composites are piecewise-linearly homotopic to the respective
identities. There are uniform primitive-recursive evaluators on rational
barycentric coordinates and moduli for the resulting maps and homotopies.
\end{theorem}
\begin{proof}
On face posets put
\[
 F(\sigma)=\{i:\sigma\subseteq S_i\},\qquad
 G(I)=\bigcap_{i\in I}S_i .
\]
The cover test makes $F(\sigma)$ nonempty, and every vertex of $\sigma$
belongs to the displayed intersection for $I=F(\sigma)$. For $I$ in
the nerve, $G(I)$ is itself a simplex of $K$: it is a nonempty subset
of any $S_i$ in the list $I$. Both assignments reverse inclusion. Sending
vertices of a chain to their images and summing the weights of repetitions
therefore defines simplicial maps between the two subdivisions. Their
composites come from monotone maps satisfying
$\sigma\subseteq GF(\sigma)$ and $I\subseteq FG(I)$.

Here is the homotopy for any monotone $u:P\to P$ with $x\le u(x)$.
Order the vertices of $P$ by a linear extension (for face posets, use size
and then mask). A point is a list of nonnegative barycentric weights
$\lambda_i$ of sum one whose positive support is a chain. Put
$A_i=\sum_{j\le i}\lambda_j$ and $A_0=0$, indexing vertices from 1.
At time $0\le t\le1$, place weight
\[
 b_i=\min(1-t,A_i)-\min(1-t,A_{i-1})
 \quad\text{at }x_i,\qquad
 \lambda_i-b_i\quad\text{at }u(x_i).
\]
Accumulate weights when vertices coincide. The bottom weights use an initial
portion of the chain and the top weights the remaining portion. Whenever
a bottom index $i$ and a top index $j$ both contribute, $i\le j$,
so $x_i\le x_j\le u(x_j)$. The top vertices form a chain by monotonicity.
The output belongs to the order complex. The formulas give the identity at
$t=0$ and $u$ at $t=1$. Zero coordinates do not affect the calculation.
This is an explicit prism homotopy, including at times when a weight splits.
With $D$ vertices the cumulative sums change by at most $D\epsilon$ for
input sup error $\epsilon$. If the time also changes by at most $\epsilon$,
each $b_i$ changes by at most $2(D+1)\epsilon$, so
$8(D+1)^2\epsilon$ bounds the change of each accumulated output. This
bound also applies to the endpoint maps on that poset.

To pass between a subdivision and the original space, send a face $A$
linearly to its barycenter. In the opposite direction the coordinate of
each nonempty face $A$ at a point $x=(x_i)$ is
\[
 w_A=|A|\max\!\left(0,\min_{i\in A} x_i-\max_{j\notin A}x_j\right),
 \qquad \max\varnothing=0 .
\]
Indeed, sort the rational coordinates $x_1\ge\cdots\ge x_v$ temporarily
(relabeling the vertices), and set $x_{v+1}=0$. Only prefixes with a
strict drop contribute. The prefix of length $j$ has weight
$j(x_j-x_{j+1})$, and its barycenter gives $x_j-x_{j+1}$ to each of
the first $j$ coordinates. Summing telescopes to the original coordinates.
Conversely, for weights on nested faces, adding their barycenters creates
one coordinate drop of exactly $w_A/|A|$ at each positively weighted
face. Thus these maps are inverse. The formula for $w_A$ is
$2v$-Lipschitz in the sup norm, and barycenter accumulation has
Lipschitz bound $2^v$. The same arguments apply to the nerve with $r$
vertices. Compose the displayed bounds, using
$D\le\max(2^v,2^r)$ for subdivisions. All formulas are finite rational
arithmetic with primitive-recursive list lengths. Evaluating input names
to precision $k+\lceil\log_2 L\rceil+1$ suffices for output precision
$k$ when $L\ge1$ is the resulting Lipschitz bound. The identities on
named points follow at each precision by these bounds. No comparison of
two exact named coordinates occurs in an evaluator.
\end{proof}

\subsection{Producing local contractions or refuting local acyclicity}

Now let $K=\bigcup_{i=1}^r K_i$ be a simplexwise cover by finite
subcomplexes. Set $K_I=\bigcap_{i\in I}K_i$ for every nonempty $I$.
Its nonemptiness is the vertex test, and $I\in N$ precisely when that
test succeeds. Chains below have coefficients in $\mathbb F_2$. The
ordinary boundary of a positive-dimensional simplex sums its codimension-one
faces, and the ordinary boundary of a vertex is zero. In the
\emph{augmented} complex add one empty face in degree $-1$ and let
every vertex have boundary that generator. In either complex
$\partial^2=0$: a pair of distinct deletions occurs in both orders. A
cycle is a vector with zero boundary. Two cycles represent the same
homology class when their sum is a boundary.

\begin{lemma}\label{lem:contraction-producer}
For a finite complex $L$, a uniform primitive-recursive algorithm either
returns a linear operator $h$ on its augmented chains, raising degree by
one and satisfying $\partial h+h\partial=1$, or returns a cycle $z$ and a
linear functional $\ell$ in the same degree with $\ell(z)=1$ and
$\ell\partial=0$ on the next degree. The first output occurs exactly
when every augmented cycle bounds.
\end{lemma}
\begin{proof}
Begin in degree $-1$ and visit basis elements $b$ by increasing degree.
With $h$ already known one degree lower, form $z=b+h\partial b$,
taking the second term as zero in degree $-1$. The identity on lower
degrees gives
$\partial z=\partial b+\partial h\partial b
=\partial b+\partial b+h\partial^2 b=0$.
Solve $\partial w=z$ over the next-degree basis, setting $h b=w$
when a solution is found. Extending linearly establishes the identity
on the whole current degree. This includes top degree, where the next
basis may be empty.

For explicit linear solving, sweep the boundary columns. Keep rows indexed
by their highest nonzero coordinate and a preimage for each row. Reduce
a new column by visiting existing pivots in descending order, adding the
pivot row whenever its lead bit is set. A nonzero remainder becomes a new
row, with preimage obtained by adding the corresponding previously kept
preimages to the new basis vector. Pivot leading coordinates are distinct.
After the sweep, reduction $R$ is linear, since each step adds the row
multiplied by one coordinate functional. Every vector is congruent to its
remainder modulo the column span, and the remainder vanishes on that span:
the rows span all input columns and reduction kills each row. Thus a zero
remainder supplies the required $w$ by the recorded preimages. Given a
nonzero remainder $Rz$, choose a set bit $j$ and output
$\ell(a)=(Ra)_j$. This bit functional annihilates each boundary column.
It proves that $z$ cannot bound. Conversely, applying the contraction
identity to any cycle gives it a filler. All loops visit listed simplices,
coordinates or pivots. This argument produces both kinds of certificate.
\end{proof}

\subsection{The double complex and naturality}

\begin{theorem}\label{thm:nerve-homology}
If Lemma~\ref{lem:contraction-producer} succeeds on each $K_I$ with
$I\in N$, then ordinary homology over $\mathbb F_2$ agrees in every degree
for $K$ and $N$. Isomorphisms in both directions can be computed on cycle
representatives. They commute on homology with simplicial maps of covered
complexes equipped with a compatible map of cover indices.
\end{theorem}
\begin{proof}
Use the finite array $D_{p,q}$ with one basis pair $(I,\sigma)$ for
$|I|=p+1$, $|\sigma|=q+1$, $\sigma\in K_I$, $p,q\ge0$.
The horizontal boundary $d_h$ deletes one index, and the vertical boundary
$d_v$ deletes one vertex, using ordinary boundaries on the two outermost
axes. Deletions on opposite axes commute, so
$d=d_h+d_v$ squares to zero. Regard a pair as having degree $p+q$.
There are two chain maps
\[
 C(K)\ \xleftarrow{\ a\ }\ (D,d)\ \xrightarrow{\ b\ }\ C(N).
\]
The first forgets a singleton $I$ and kills the terms with $p>0$. The
second forgets a singleton $\sigma$ and kills terms with $q>0$.
For $a$, the two ways to delete from a two-element $I$ give the same
$\sigma$ and cancel after forgetting. Other horizontal terms vanish after
$a$, and $a$ commutes with vertical deletions. The same computation with
axes exchanged proves the assertion for $b$.

Augment $D$ on the left by the column $D_{-1,q}=C_q(K)$, letting
horizontal deletion of a singleton act as $a$. For fixed $\sigma$ the
horizontal complex is the augmented simplex on the nonempty set
$T_\sigma=\{i:\sigma\in K_i\}$. Put $c=\min T_\sigma$. A horizontal
contraction sends $I$ to $I\cup\{c\}$ if $c\notin I$, and to zero
otherwise, also sending the empty face to $\{c\}$. If $c\notin I$,
deleting $c$ returns $I$, while all other deletions pair with the coned
faces of $I$. If $c\in I$, only deletion of $c$ contributes to coning
the boundary. This proves the contraction identity in both cases.

Any cycle $u$ in this augmented array has a filler. Sweep the rows by
decreasing $q$. In the current highest remaining row, the cycle equation
says $d_h u_q=0$, as there is no incoming term from a higher row of $u$.
Apply the horizontal contraction to $u_q$, obtaining $v_q$ with
$d_hv_q=u_q$. Replace $u$ by $u+d v_q$. Its top row vanishes, its
cycle condition persists, and all new terms have smaller $q$. At $q=0$
there is no outgoing vertical boundary. The sum of recorded $v_q$ fills
the original cycle. Empty rows can be visited with zero filler, so the
sweep length is fixed in advance.

Independently augment downward by $D_{p,-1}=C_p(N)$ using $b$. For
each $I$, the vertical complex is exactly the augmented complex of $K_I$.
Use the $h_I$ from Lemma~\ref{lem:contraction-producer} and the identical
sweep in decreasing $p$. This fills all cycles in the second augmented
array.

To see what these fillers compute, write an augmented chain as $(x,y)$,
where $x$ is in $D$ in degree $k$ and $y$ is in the target in degree
$k+1$. Its boundary is $(dx,f x+\partial y)$ for $f=a$ or $f=b$.
For a target $0$-cycle this convention puts $(0,y)$ in total degree
$-1$; all ordinary chain groups below degree zero are zero. If $K$ is
empty, so is $N$, and both ordinary homology groups vanish in every
degree. Given a target cycle $y$, fill $(0,y)$ to get $(z,w)$. Then $dz=0$
and $fz+\partial w=y$, which lifts the class. Given a cycle $x$ whose
image bounds, say $fx=\partial y$, fill $(x,y)$ to get a $z$ with
$dz=x$. Thus each map reflects boundaries and induces an isomorphism.
Applying one projection after lifting through the other gives the two
claimed algorithms. Choices of lifts give homologous results by boundary
reflection.

For naturality let $g$ send vertices simplicially to another covered
complex and $s$ send cover indices with $g(K_i)\subseteq K'_{s(i)}$.
Use $(I,\sigma)\mapsto(s(I),g(\sigma))$ when both maps are injective
on those faces, and zero otherwise. This respects the incidence condition.
For one axis, deleting commutes with the mapping rule when there is no
duplicate image. With just one duplicate pair, the two deletions that
could remove the duplication have the same image and cancel. With any
further duplication every deletion still has a duplicate. This proves
the chain-map rule on both axes and also for the maps of $K$ and $N$.
Both projections commute with these maps, by inspection on singleton
faces and on faces of larger size. Invert the two projection isomorphisms
on homology to get the naturality square. In particular the comparison
commutes with every supplied refinement in a tower of such covers.
\end{proof}

\paragraph{Bounds and interpretation.}
Let $V=2^v$ and $R=2^r$. There are fewer than $R$ intersections, each
with at most $V$ augmented basis vectors. Each pivot calculation visits
at most $V$ columns with at most $V$ pivot additions per reduction, on
vectors of at most $V$ bits. Evaluating the obstruction functional on
the original basis costs at most $V$ further reductions in the failing
degree. Boundary and contraction matrices can be stored in $V^2$ bits
per matrix, and the augmented comparison arrays have at most $RV$ entries
as chain masks. The two sweeps use at most $v$ rows or $r$ columns;
all nested loops therefore have specified primitive-recursive limits.
Finite bitwise identities such as $\partial h+h\partial=1$ can be checked
on all basis columns by a bounded Boolean fold. For the sweep the
inductive assertion is that all processed rows vanish, the remainder is
closed, and its sum with the boundary of the accumulated filler is the
original input. Each clause is a bounded equality of arrays at the fixed
input size. In first-order iPRA one can state boundary reflection as the
universal closure of a conditional between such equalities, supplying the
input cycle and a purported image filler as numbers. A failed local test
certifies failure of the acyclicity hypothesis. Some particular nerves
can still have the right homology by other arguments.

\paragraph{Example.}
Take the boundary of a hexagon with vertices $0,\ldots,5$. Three subcomplexes
are the two-edge paths $0,1,2$, $2,3,4$, and $4,5,0$. Each pair has a
single common endpoint, and the triple intersection is empty. The nerve
has three edges and no face of dimension two. The accompanying implementation
lifts the full hexagon cycle and projects it to the three-edge cycle;
the reverse transfer gives the original hexagon. Covering the same hexagon
with just one member instead produces the six-edge obstruction cycle.
The dual functional tests a single edge. It evaluates to one on the
cycle and vanishes on the boundaries from degree two because there are no
two-simplices. An independent small-input check exhausts all $19^2$
ordered pairs of complexes on three labeled vertices as covers of their
union. Literal enumeration of cycles and boundaries identifies 109
admissible covers and 252 local obstructions, agreeing with the solver.
It also checks 1,553 cycle round trips modulo boundaries and 1,788
naturality squares for index permutations and admissible index mergers.

\paragraph{Lean completion path.}
The companion \texttt{NervePolarity.lean} constructs an actual closed
\texttt{Prf PRAAx} derivation for arbitrary unary and binary arithmetic formulas $A,R$ of
$\forall x\,i\,[(A(x)\land\forall z(A(z)\to R(z,i)))\to R(x,i)]$.
With $R$ the incidence relation, this is the logical unit $A\subseteq GF(A)$;
transposing incidence supplies the dual unit. The file
\texttt{NerveTransfer.lean} checks
\texttt{bounded\_eliminate}, the algebraic filtration argument, and
\texttt{lift\_cycle} and \texttt{reflect\_boundary}, the two augmented-complex
inferences. The latter expose all hypotheses in \texttt{ConeExact}.
To instantiate them, define a simplex as a bounded mask and a chain as its
Boolean coefficient array. Construct deletion by toggling the indicated
bit and define each differential by a parity fold. Pair the two deletion
orders to prove $d^2=0$. Implement the pivot scan with its leading-coordinate
and preimage invariants, then store $h$ as a table of vectors. Instantiate
the stage predicate of \texttt{bounded\_eliminate} by vanishing above a
row (or column) in the augmented array. The explicit local operators in
the proof supply its step and the two projections give its
\texttt{ConeExact} instances. In the object language expand each of these
finite operations at natural-number type and use the indicated Boolean
invariants in quantifier-free induction. For Theorem~\ref{thm:nerve-pl}
the independent implementation path consists of face inclusion, incidence
transpose, ordered cumulative sums, and the min/max formulas. The order
inequalities and coordinate identities are rational tests on finite
arrays. Denominator clearing gives their object-language form. A semantic
Lean proof of these operations and an object derivation of their identities
are separate outstanding steps.

\subsection{Producing a homotopy inverse from collapsible intersections}
\label{sec:collapse-nerve}

The linear contraction in Lemma~\ref{lem:contraction-producer} fills
chains. A geometric contraction allows one to fill maps, which gives a
stronger conclusion for a decidable class of covers. Use the subcomplex
cover $K=\bigcup_i K_i$, intersections $K_I$, and nerve $N$ of
Section~\ref{sec:nerve}. A legal elementary collapse removes a pair
$(\sigma,\tau)$, where $\sigma$ is nonempty, $\tau=\sigma\cup\{v\}$
for one new vertex $v$, and $\tau$ is the only face strictly containing
$\sigma$ in the current complex. Call a complex collapsible if a finite
sequence of such removals leaves a single vertex.

\begin{theorem}\label{thm:collapse-nerve}
There is a primitive-recursive test whether every nonempty $K_I$ is
collapsible. On success it produces rational evaluators with PR moduli
for maps $|K|\rightleftarrows |N|$ and homotopies from both composites to
their respective identities. On failure it identifies an intersection
without a collapse sequence of the indicated type.
\end{theorem}
\begin{proof}
\emph{Local production.}
If an intersection has $S$ nonempty faces, each move removes two.
Enumerate words of length at most $S$ over at most $S^2$ possible
pairs, check each move in its current face table, and test the final
vertex condition. The number of words is at most $(1+S^2)^{S+1}$.
This search includes every successful collapse sequence. Perform it on
each of the at most $2^r-1$ nonempty index sets, omitting empty
intersections. An unsuccessful search excludes collapsibility for this
triangulation; a different triangulation or a different contraction
can have different certificates.

One elementary move supplies an actual deformation to the remaining
subcomplex. For barycentric coordinates $x$ in the current complex put
$m(x)=\min_{j\in\sigma}x_j$. At time $t$, subtract $t m(x)$ from each
coordinate in $\sigma$ and add $|\sigma|t m(x)$ to the coordinate $v$.
If $m(x)>0$, the support contains $\sigma$, so it lies in $\tau$.
The altered point stays in that simplex, with nonnegative coordinates
of sum one. At $t=1$ at least one coordinate in $\sigma$ vanishes. If
the input lies in the remaining subcomplex its support misses a vertex
of $\sigma$, so the entire motion fixes it. These assertions also cover
the case where $\sigma$ is a vertex. In the sup norm on both coordinates
and time, the move has Lipschitz bound $2|\sigma|+1$: the minimum changes
by at most the coordinate error, and multiplication by $t\in[0,1]$
adds at most the time error.

If a sequence has length $s$, apply its moves in order at respective
times $t_j=\min(1,\max(0,st-j))$, for $0\le j<s$.
All preceding moves are complete whenever a move is active. Inactive
moves have coefficient zero in the same coordinate formula. This
constructs a rational contraction $H_I(x,t)$ from identity to its
terminal vertex. For input on at most $v_0$ vertices, a common
Lipschitz bound in the product sup norm is obtained by $B_0=1$,
$B_{j+1}=(2v_0+1)(B_j+s+1)$, taking $B_s$ at the end. A zero-length
sequence occurs on a single vertex and uses the identity there.

\emph{Extension of maps in nested carriers.}
Here a carrier for a simplex $e$ of a finite complex $L$ is a
subcomplex $C(e)$ of the target with $C(e')\subseteq C(e)$ for every
nonempty face $e'$ of $e$. Suppose each $C(e)$ comes with a rational
contraction $H_e$ from identity to a specified point and with a Lipschitz
bound. A map already defined on a subcomplex, with restrictions Lipschitz
and carried simplexwise, extends constructively over $L$. We describe
the construction because it supplies the uniformity needed in iPRA.

Initialize missing vertices at the contraction endpoints. For a missing
simplex with $n\ge2$ vertices the entire boundary has already been mapped
into $C(e)$. Write its map as $b$ and suppose its restrictions to facets
have Lipschitz bound $L_0\ge1$. The whole boundary has bound $2nL_0$
in barycentric sup norm. Indeed, if rational points $x,y$ have distance
$\delta$ with $n\delta\ge1$, all target coordinates lie in $[0,1]$.
Otherwise their supports meet: the mass outside the common support is at
most $n\delta$. Form $z$ by taking coordinatewise minima and moving
the remaining mass, at most $n\delta$, to any common vertex.
Both $x,z$ lie in one source facet and $y,z$ in another;
$\|x-z\|_\infty,\|y-z\|_\infty\le n\delta$.
Agreement of the facet maps on their common face yields the bound.
The same argument glues simplexwise Lipschitz bounds over any finite
source complex, using its total number of vertices.

Write $c=(1/n,\ldots,1/n)$ and
$u(x)=\max_i(1-n x_i)$. It ranges from $0$ at $c$ to $1$ on the
boundary. For $u>1/2$ put $y=c+(x-c)/u$; the choice of $u$ ensures
that the coordinates of $y$ are nonnegative, sum to one, and include
a zero. Extend by
\[
  E(x)=
  \begin{cases}
    H_e(b(y),2(1-u)),&u>1/2,\\
    H_e(\text{endpoint},1),&u\le1/2.
  \end{cases}
\]
This agrees with $b$ at $u=1$ and is constant on the inner half-sized
simplex. For two outer inputs of distance $\delta$, their radial
parameters differ by at most $n\delta$, their inverses by at most
$4n\delta$, and the boundary arguments differ by at most
$(2+4n)\delta$. At an inner/outer pair, compare the outer expression
with $H_e(b(y),1)$ using the same outer boundary argument; the time
difference is at most $2n\delta$. Thus, if $B$ bounds the contraction,
\[
 B\bigl[2nL_0(2+4n)+2n+1\bigr]
\]
bounds the extension in all cases. At a rational boundary argument one
may evaluate any least indexed facet containing it, since the facet
maps agree there. The threshold comparison takes place on rational
samples. The displayed modulus extends the result to names. Visiting
simplices in increasing dimension proves the extension assertion.

The same construction gives a carried homotopy between two carried maps,
with known Lipschitz bounds, by extending their prescribed values at
the two ends of a cylinder. To make the cylinder data explicit, order
the vertices of $L$ once. For a simplex on $x_0<\cdots<x_d$ triangulate
its cylinder with the simplices
\[
 (x_0,0),\ldots,(x_j,0),(x_j,1),\ldots,(x_d,1)
 \qquad(0\le j\le d)
\]
and their faces. The assignment restricts compatibly on faces. The
carrier of a cylinder simplex is $C(e)$ for its projected source face
$e$. The identification of the product with this triangulation splits
the ordered barycentric masses at total mass $1-t$: assign the initial
mass to time-zero vertices and the remainder to time-one vertices,
using the $b_i$ formula in Theorem~\ref{thm:nerve-pl}. Its support
is contained in one displayed simplex. The inverse adds the two
coordinates at each source vertex and takes the total top mass as time.
Both operations have explicit Lipschitz bounds by finite summation and
minimum. Thus the extension produces a homotopy with the stated
endpoints also across every shared face.

\emph{The two comparison maps.}
For a face $\sigma$ of $K$, put $F(\sigma)=\{i:\sigma\in K_i\}$.
This is a nerve face and $F$ reverses inclusion, giving a simplicial
map $F:|\operatorname{sd}K|\to|\operatorname{sd}N|$ by accumulating
weights. Identify subdivisions with their originals by the rational
coordinate formulas of Theorem~\ref{thm:nerve-pl}. To construct
$G:|\operatorname{sd}N|\to |K|$, give a chain
$e=(I_0\subsetneq\cdots\subsetneq I_d)$ the carrier $K_{I_0}$.
Removing its initial vertex can only shrink the carrier. Apply the
extension just proved, starting with its chosen vertices. In particular
$G(e)\subseteq |K_{I_0}|$, with $G(e)$ here denoting the image of
the whole geometric simplex.

For the composite $GF$, a source chain
$\sigma_0\subsetneq\cdots\subsetneq\sigma_d$ has carrier
$K_{F(\sigma_d)}$. Both its identity image inside $\sigma_d$ and the
image under $GF$ lie there: all $F(\sigma_j)$ contain $F(\sigma_d)$,
and the image of the entire chain is a face of the reverse chain with
least element $F(\sigma_d)$. These carriers are nested on faces and
have the produced contractions. The cylinder construction gives the
first homotopy.

For $FG$ use the order complex of the poset of all nerve faces
containing $I_0$ as carrier for $e$. It contracts to the vertex $I_0$
by $(1-t)x+t\mathbf e_{I_0}$: adjoining the least element to any chain
still gives a chain. Its Lipschitz bound is at most three. The identity
is carried there. Every point of $|K_{I_0}|$, when expressed in
subdivision coordinates, has its positively weighted faces $\sigma$
in $K_{I_0}$. For rational points this follows from the explicit
coordinate formula: each such $\sigma$ is a subset of the point's
support. Thus $F(\sigma)\supseteq I_0$, and $FG$ is carried there too.
Apply the cylinder construction again. All verifications on rational
points pass to names with the prescribed moduli and the resulting
uniform equality estimates.

Finally, one may bound the number of subdivision vertices by
$D=\max(1,2^{v_0},2^r)$ and the number of simplices needed in any
one cylinder by $2^{2D}$. At each inductive stage take the maximum
of the preceding bound and the displayed extension bound over the
finite list, then proceed to the next dimension. A rational point
evaluation follows at most one descending face chain per extension;
each step uses a bounded facet scan, rational arithmetic and a
contraction with a previously bounded number of moves. The successive
constructions of $G$, the endpoint composites and the two homotopies
have fixed nesting depth as constructions. Table enumeration, the
dimension recurrence for the Lipschitz constants and these evaluation
clocks are PR. An output tolerance $2^{-k}$ is obtained by evaluating
a rational input sample to tolerance $2^{-k}/(2L)$ for the resulting
bound $L\ge1$. Empty realizations use the empty maps. This proves
the claim for the represented spaces.
\end{proof}

In the hexagon example each two-edge arc collapses by removing its
two outer vertices with their incident edges, leaving its middle vertex.
Each nonempty pairwise intersection already consists of one vertex.
The theorem therefore upgrades the computed homology transfer for that
cover to a homotopy equivalence with the triangular nerve.

For Lean, add to the chain-comparison targets three independent
finite kernels: the free-face test and min-coordinate collapse, the
collar extension with its two cases, and the cylinder incidence list.
Record the extension invariant as agreement on previously processed
faces, membership in the assigned carrier, and the displayed Lipschitz
estimate. For an object proof, keep the two input sample points and the
requested accuracy as parameters. For each fixed construction, carry the
invariant through the descending evaluator path and its bounded table
lookups; a finite trace records all rationals and comparisons used in
both evaluations. Equivalently, bound their codes in advance by iterating
a majorant for one rational operation through the PR operation clock,
then Booleanize the invariant for all such subcalls at each depth.
The geometric maps and homotopies in this section are proved on paper;
the supplied \texttt{ElementaryCollapse.lean} checks mass preservation, nonnegativity, the free-face endpoint, and the fixed-wall property of the local deformation. The carrier extensions and their assembly remain deep formalization targets.

\subsection{A one-member cover separating the two tests}
One can see the precise difference between geometric collapse and linear filling on 13 vertices. Take $a_i$ for $i$ modulo $9$, $b_j$ for $j$ modulo $3$, and a vertex $p$. For each $i$, put $j=i\bmod3$ and include the three triangles
\[
 (p,a_i,a_{i+1}),\qquad (a_i,a_{i+1},b_{j+1}),\qquad
 (a_i,b_j,b_{j+1}),
\]
with all their faces. Cover the resulting complex $K$ by $K$ alone. Its nerve is a single vertex. There are 13 vertices, 39 edges and 27 triangles. We compute
\[
 (\dim H_0,\dim H_1,\dim H_2)(K;\mathbb F_2)=(1,0,0),
 \qquad
 (\dim H_0,\dim H_1,\dim H_2)(K;\mathbb F_3)=(1,1,1).
\]
To see this, temporarily replace each $b_{i\bmod3}$ at the outer rim by a separate vertex $b'_i$. The fan from $p$ fills the inner nine-gon, and the other 18 triangles triangulate its annular collar. Orient the resulting disk consistently. Along any internal edge the two triangle coefficients of a cycle must agree, as the induced edge signs are opposite. The triangles are connected by internal edge adjacencies, so all coefficients are a common scalar $c$. Upon restoring the identifications of rim vertices, only rim edges are identified: each edge $b_jb_{j+1}$ receives the same oriented rim contribution three times. Thus $\partial_2$ kills the constant coefficient vector exactly when $3c=0$, and these exhaust its kernel. The connected edge graph gives $\operatorname{rank}\partial_1=12$ (a spanning tree gives 12 independent differences, and every edge difference has coordinate sum zero). Hence $\dim\ker\partial_1=27$; subtracting the rank of $\partial_2$ proves the formulas.

The one-member cover passes the local $\mathbb F_2$ contraction test. It has no elementary collapse to a point: every edge has at least two incident triangles, and every vertex lies in several triangles. Its $\mathbb F_3$ homology also separates it from a point. In particular the $\mathbb F_2$ test alone cannot establish the conclusion of the geometric nerve theorem; even its predicted agreement of homology with other coefficients can fail. All displayed ranks are computations on finite signed boundary matrices. The companion \texttt{threefold.py} computes them independently by modular elimination.

